\documentclass[11pt,a4paper]{article}
\usepackage[utf8]{inputenc}
\usepackage[spanish]{babel}
\usepackage{geometry}
\usepackage{booktabs}
\usepackage{tabularx}
\usepackage{xcolor}
\usepackage{hyperref}

\geometry{top=2.5cm, bottom=2.5cm, left=2.5cm, right=2.5cm}
\definecolor{brandgreen}{RGB}{27, 77, 62}

\title{\textbf{\color{brandgreen} PDAEH -- Plantilla de Artefacto: CPD} \\[0.3em] \large Catálogo de patrones de diseño}
\author{Proceso para la Definición de Arquitecturas Empresariales Híbridas}
\date{\today}

\begin{document}
\maketitle

\noindent \textit{\textbf{Nota sobre el uso de esta plantilla:} Esta plantilla es una guía metodológica sugerida y no constituye un instrumento rígido o mandatorio. Cada organización, empresa, grupo de investigación o equipo de trabajo es libre de adaptarla, extenderla o simplificarla según sus políticas internas, herramientas y nivel de madurez.}

\section*{1. Propósito y Alcance del Artefacto}
Orientar la resolución técnica recurrente a nivel de interacción híbrida mediante patrones de diseño estandarizados.

\section*{2. Estructura y Registro del Contenido}
\begin{tabularx}{\textwidth}{|>{\bfseries}p{4.5cm}|X|}
    \hline
    \rowcolor[gray]{0.92} Campo / Sección & Registro del Contenido \\ \hline
    \textbf{Nombre} & [Completar información] \\ \hline
    \textbf{Problema y contexto} & [Completar información] \\ \hline
    \textbf{Fuerzas} & [Completar información] \\ \hline
    \textbf{Solución} & [Completar información] \\ \hline
    \textbf{Participantes} & [Completar información] \\ \hline
    \textbf{Interacciones} & [Completar información] \\ \hline
    \textbf{Consecuencias} & [Completar información] \\ \hline
    \textbf{Ejemplo híbrido} & [Completar información] \\ \hline
\end{tabularx}

\vspace{1.2cm}
\section*{3. Control de Revisiones y Gobierno Arquitectónico}
\noindent \textbf{Elaborado por:} \hrulefill \\[0.5em]
\textbf{Fecha:} \hrulefill \\[0.5em]
\textbf{Aprobación del Órgano de Gobierno Arquitectónico (OGA):} Aprobado [ \, ] \quad Rechazado [ \, ]

\end{document}
